
\documentclass[11pt,a4paper]{article}

\usepackage[margin=1in]{geometry}
\usepackage{amsmath}
\usepackage{enumitem}

\title{HW01 -- Robotics Assignment Report}
\author{Tomer Sananes & Gil Pasi}
\date{}

\begin{document}

\maketitle

\section{Assignment Objective}

The objective of this assignment is to detect a red marker using a camera, calculate its position in 3D space, and move a robotic arm to a target point above the marker without touching it.

\section{Solution Overview}

\subsection{Step 1 -- Receiving Camera Data}

The system receives an RGB image, a depth image, and camera calibration information (CameraInfo).

The RGB and depth images are synchronized using their timestamps to ensure that both measurements correspond approximately to the same moment.

\subsection{Step 2 -- Marker Detection}

The RGB image is converted to the HSV color space to simplify red color detection.

Noise filtering is then applied, and the largest connected red region is selected. The center of this region is used to estimate the marker's pixel position.

The marker's depth is calculated using the median of the valid depth values within the detected region, reducing the influence of outliers.

\subsection{Step 3 -- Calculating the 3D Position}

Using the \textbf{Pinhole Camera Model}, the marker's pixel coordinates, depth, and camera calibration parameters are used to calculate its position in the camera coordinate frame:

\begin{equation}
    X = \frac{(u-c_x)Z}{f_x},
    \qquad
    Y = \frac{(v-c_y)Z}{f_y}
\end{equation}

The marker's position is then transformed into the robot's base coordinate frame (\texttt{base\_link}) using the corresponding transformation.

Finally, 8 cm is added to the Z coordinate. This value represents the \textbf{fixed height difference between the marker and the target position}, as specified in the assignment. The resulting point is used as the robotic arm's movement target.

\section{Code Structure}

\begin{itemize}[leftmargin=*]
    \item \textbf{student\_solution.py} -- The main file, responsible for receiving camera data, synchronizing images, publishing the marker's position, and sending movement commands to the robotic arm.

    \item \textbf{utils.py} -- Contains helper functions for image processing, marker detection, 3D position calculation, and coordinate transformations.
\end{itemize}

\section{Summary}

We implemented a system that detects a red marker from camera images, calculates its 3D position, and automatically moves the robotic arm to the target point above it.

\end{document}
